\section{Introduction}
\label{sec:introduction}

% 引言是开展研究的动机，让评论员觉得这个工作该做。详细的叙述放到相关工作

% 第二段，冗余分析（事先），因为浪费很多，所以需要观测规划
% 第三段，根据第二段引出在线的测绘规划的必要性
% 作为我的方法的特点，（闭环的体系，有规划结果）（通用性）

For missions to explore distant small astroids(NASA's OSIRIS-REx\citep{alasad2021}, ESA's Rosetta\citep{taylor2017rosetta}, JAXA's Hayabusa\citep{fujiwara2006,demura2006}, etc), reliable three-dimensional models is essential throughout multiple scientific tasks. Stereophotoclinometry (SPC)\citep{barnouin2020} is an established approach to extract information from captured obtical images and thus derive models of astroids. It connects pixels of multiple images and illumination status to stereo constraints like elevation and albedo of a specific surface shape, making which integratable variables. Its practical use across planetary missions demonstrates its value of extrating data in roughly unknow environments with limited informations. However, the image collections of real missions for SPC reconstruction are massive. Rosetta's SHAP5 model of comet 67P incorporated more than 5500 narrow-angle camera images acquired through June 2015 \citep{preusker2017}. For OSIRIS-REx, successive Bennu models v10, v20, and v42 used 2351, 5114, and 15,325 images\citep{palmer2022}. To obtain reliable SPC inputs, observation plan incorporated repeat passes over target astroid with overlapping imaging sequences, creating substantial redundancy. This redundancy provides margin for operation teams, but massive images with similar observation geometries can yield diminishing improvements while continuing to consume observation time, image storage, downlink capacity, and human processing effort. A planning simulations in commet 67P scenarioal so indicates avoidable over-observation by 43\%, while improving the geometric mapping score \citep{piccinin2022}. Therefore, a well-planned observation strategy is essential to ensure that the images acquired for SPC reconstruction are both scientifically valuable and operationally efficient.

% Acquiring additional images without assessing their contribution can therefore lead to inefficient image use and wasted observing resources, motivating a planner that directs further observations toward unresolved terrain and stops when the required model quality has been achieved.

% Because the surface topography and albedo are totally or partially unknown before reconstruction, the observation geometry is the only controllable factor that can be used to improve the quality of reconstruction, therefore the quality of modeling depends on both the reconstruction procedure and the observations made available to it. The brightness of a pixel on images depens on not only the surface topography but also the observation geometry(illumination conditions and camera extrinsic parameters). Stereo observations of multiple images constrain the three-dimensional positions of the surface, while photoclinometry estimates local slopes and relative albedo from image brightness; these slopes are integrated into surface elevations with geometric and topographic constraints. Repeatedly imaging a specific area with the same illumination condition may therefore increase its nominal coverage without adequately constraining its elevation\citep{gaskell2023,palmer2022}, which makes observation of low infomation density. This limitation makes the choice of images a scientific determinant of model quality: an observation is valuable insofar as it adds constraints that the existing image set lacks. All above makes observation planning, the outcome of a sequence of operational decisions, an important scientific task for SPC model reconstruction.

Existing works generate observing plan of astroids by optimizing the SPC-inspired observation geometry and coverage scores\citep{pesce2018mapping}, which later be used in deep reinforcement learning to coordinate orbit selection and imaging \citep{chan2019mapping}. \citet{piccinin2022} learned when to acquire images along a prescribed trajectory, reducing image count while improving geometric mapping scores and transferring the policy across targets. \citet{wang2025} combined a genetic algorithm with local search to reduce observing positions and images under illumination, visibility, and coverage constraints. However, while SPC is an algorithm which generates and refines models round after round, these planners don't take regional quality diagnostics of a newly reconstructed SPC model from last round into consideration, and their reported evaluations don't bennefit the final SPC model, they omit the modeling process and its feedback. Such scores describe the expected suitability of the input images rather than the accuracy of the real model, sometimes favorable aggregate values can obscure poorly constrained regions. Reconstruction-based studies instead evaluate observation choices through their modeling outcomes. \citet{palmer2024geometry} reconstructed synthetic terrain under different imaging configurations and derived complementary viewing and illumination recommendations for mission planning. For existing image archives, \citet{zhang2025imagepairs} used KD-tree searches to accelerate overlap analysis and select SPC inputs, then generated a terrain model and quantified elevation errors against a reference model; this addresses image selection rather than the planning of new acquisitions. To make subsequent acquisitions respond to actual modeling needs, an automated planner therefore requires regional SPC feedback that identifies unresolved terrain and updates observation priorities as the reconstruction improves. Producing and validating final shape models or terrain products is also necessary to demonstrate that planning gains translate into the required accuracy, using known truth in simulations or independent measurements where available in flight, because internal consistency alone can miss shared systematic errors \citep{alasad2021,palmer2024accuracy}.

Complementing observation-planning research, SPC validation and sensitivity studies have assessed the geometric accuracy of reconstructed models and the contributions of their input images. \citet{alasad2021} developed a validation framework for Bennu shape models that combined stereo residuals, height dispersion among overlapping maplets, comparisons between rendered and observed images, and independent laser-altimeter measurements, enabling model quality to be assessed beyond internal consistency. \citet{weirich2022} rendered images of a synthetic asteroid with known topography, reconstructed its global shape using SPC, and compared the reconstruction with the truth model after registration to calibrate internal quality metrics against measured 3D error. \citet{palmer2024accuracy} varied image sets and processing conditions in synthetic-terrain experiments, combining terrain differences with normalized cross-correlation between simulated observations and images rendered from reconstructed models to evaluate geometric error and reveal deficiencies in fine surface detail. Focusing on the role of imaging conditions, \citet{palmer2024geometry} compared SPC reconstructions from prescribed combinations of viewing and illumination geometries and image counts, identifying a compact set of four complementary topographic images and one albedo image near local noon that produced high-quality terrain in their tests. These studies provide quantitative methods for assessing reconstruction accuracy and evidence of how image combinations support terrain and albedo recovery. These evaluations assess already acquired or simulated image sets after SPC reconstruction, explaining the contribution of the selected images through the resulting model quality. Although their findings can inform observation design, they do not provide an acquisition planner that predicts a candidate image's accuracy gain before capture and uses the latest regional reconstruction diagnostics to select subsequent observations.

To address these limitations, this paper develops an observation planning algorithm that directly integrates SPC reconstruction and accuracy assessment into a deep reinforcement learning framework. The decision model uses the reconstructed surface and regional quality diagnostics from the previous round to select observation positions for the next round, forming a closed loop of observation planning, SPC reconstruction, accuracy assessment, and replanning. Compared with manually planned observation campaigns, the proposed method achieves comparable reconstruction accuracy while acquiring fewer images, thereby conserving imaging resources. Compared with existing automated SPC observation planners evaluated primarily through observation geometry and coverage scores, the proposed framework delivers final SPC surface models as data products and quantitatively assesses their geometric accuracy, enabling planning performance to be evaluated in terms of both imaging efficiency and the quality of the resulting models.

